\documentclass[a4paper,12pt]{article}

\usepackage[T2A]{fontenc}
\usepackage[utf8]{inputenc}
\usepackage[russian]{babel}

\usepackage{amsmath,amssymb,amsthm,mathtools}
\usepackage{helvet}
\renewcommand{\familydefault}{\sfdefault}
\usepackage{icomma}

\usepackage[
  a4paper,
  left=20mm,
  right=20mm,
  top=20mm,
  bottom=22mm
]{geometry}

\usepackage{microtype}
\usepackage{tikz}
\usepackage{tkz-euclide}
\usepackage{pgfplots}
\pgfplotsset{compat=1.18}

\usepackage{tabularx,booktabs,longtable}
\usepackage{enumitem}
\usepackage{hyperref}
\usepackage{parskip}
\usepackage{fancyhdr}
\usepackage{setspace}
\usepackage{xcolor}

\definecolor{accent}{HTML}{008080}
\definecolor{darkaccent}{HTML}{004D4D}
\definecolor{textgray}{HTML}{333333}
\definecolor{softgray}{HTML}{F2F4F4}

\providecommand{\tg}{\mathop{\rm tg}\nolimits}
\providecommand{\ctg}{\mathop{\rm ctg}\nolimits}
\providecommand{\arctg}{\mathop{\rm arctg}\nolimits}

\hypersetup{
  colorlinks=true,
  linkcolor=darkaccent,
  urlcolor=darkaccent
}

\color{textgray}
\onehalfspacing
\setlength{\parindent}{0pt}
\setlength{\emergencystretch}{3em}
\raggedbottom

\setlist[itemize]{
  leftmargin=1.6em,
  itemsep=0.28em,
  topsep=0.35em
}

\setlist[enumerate]{
  leftmargin=*,
  itemsep=0.7em,
  topsep=0.4em
}

\newcommand{\checkitem}{\item[\textcolor{accent}{$\square$}]}
\newcommand{\key}[1]{\textcolor{darkaccent}{\textbf{#1}}}

\newcommand{\sectionrule}{%
  \par\vspace{-0.35em}%
  {\color{accent!60}\rule{\linewidth}{0.7pt}}%
  \vspace{0.35em}
}

\pagestyle{fancy}
\fancyhf{}
\fancyhead[L]{\small\textcolor{textgray}{Прямоугольный треугольник}}
\fancyhead[R]{\small\textcolor{textgray}{София Кальней}}
\fancyfoot[C]{\small\thepage}
\renewcommand{\headrulewidth}{0.4pt}
\renewcommand{\headrule}{\hbox to\headwidth{%
  \color{accent!55}\leaders\hrule height \headrulewidth\hfill}}

\tkzSetUpLine[line width=1pt,color=darkaccent]
\tkzSetUpPoint[color=darkaccent,fill=darkaccent,size=3]
\tkzSetUpLabel[color=black,font=\small]

\begin{document}

% =========================================================
% TITLE PAGE
% =========================================================
\begin{titlepage}
\thispagestyle{empty}

\vspace*{2.5cm}

\begin{center}

{\Large\textcolor{accent}{ОГЭ · геометрия}\par}

\vspace{1.1cm}

{\Huge\bfseries Чек-лист\par}

\vspace{0.6cm}

{\LARGE\bfseries
Прямоугольный треугольник:\par}

\vspace{0.15cm}

{\LARGE
медиана, высота, площадь\\
и тригонометрия\par}

\vspace{1.3cm}

{\color{accent!70}\rule{0.58\linewidth}{1pt}}

\vspace{1.3cm}

{\large
Лёвин Артём Александрович\\[0.25em]
\textbf{эксклюзивно для Софии Кальней}
\par}

\vspace{0.9cm}

{\large \today\par}

\vfill

{\small\textcolor{textgray}{
Ключевые свойства, алгоритмы и типовые конструкции
для решения геометрических задач
}\par}

\end{center}

\vfill

\begin{tikzpicture}[remember picture,overlay]
\draw[
  accent!55,
  line width=1.1pt
]
([xshift=1.5cm,yshift=1.45cm]current page.south west)
.. controls
([xshift=5.0cm,yshift=2.05cm]current page.south west)
and
([xshift=-5.0cm,yshift=0.95cm]current page.south east)
..
([xshift=-1.5cm,yshift=1.45cm]current page.south east);
\end{tikzpicture}

\end{titlepage}

% =========================================================
% CHECKLIST
% =========================================================
\section*{\textcolor{darkaccent}{1. Базовая конструкция: медиана в прямоугольном треугольнике}}
\sectionrule

Пусть в треугольнике \(ABC\)
\[
\angle ACB=90^\circ,
\]
а \(CM\) --- медиана, проведённая к гипотенузе \(AB\).

\begin{itemize}
  \checkitem Медиана делит гипотенузу пополам:
  \[
  AM=MB=\frac{AB}{2}.
  \]

  \checkitem В прямоугольном треугольнике середина гипотенузы равноудалена
  от всех трёх вершин:
  \[
  \boxed{AM=BM=CM=\frac{AB}{2}}.
  \]

  \checkitem Следовательно, точка \(M\) является центром описанной окружности,
  а радиус этой окружности равен
  \[
  R=\frac{AB}{2}.
  \]

  \checkitem Треугольники \(AMC\) и \(BMC\) равнобедренные:
  \[
  AM=CM,\qquad BM=CM.
  \]
\end{itemize}

\begin{center}
\begin{tikzpicture}[scale=0.9]
  \tkzDefPoint(0,0){C}
  \tkzDefPoint(0,3){A}
  \tkzDefPoint(4,0){B}
  \tkzDefMidPoint(A,B)
  \tkzGetPoint{M}

  \tkzDrawSegments(A,B B,C C,A C,M)
  \tkzMarkRightAngle[size=0.28](A,C,B)

  \tkzLabelPoints[above](A)
  \tkzLabelPoints[right](B)
  \tkzLabelPoints[below left](C)
  \tkzLabelPoints[above right](M)
\end{tikzpicture}
\end{center}

\key{Обратное утверждение.}
Если в треугольнике медиана, проведённая к стороне \(AB\), удовлетворяет
\[
CM=\frac{AB}{2},
\]
то
\[
\angle ACB=90^\circ.
\]

Иными словами, равенство
\[
CM=AM=BM
\]
позволяет распознать прямоугольный треугольник.

% =========================================================
\section*{\textcolor{darkaccent}{2. Диаметр и прямой угол}}
\sectionrule

Если вершина \(C\) лежит на окружности, а \(AB\) является диаметром, то
\[
\boxed{\angle ACB=90^\circ}.
\]

Это теорема о вписанном угле, опирающемся на диаметр.

\begin{center}
\begin{tikzpicture}[scale=0.8]
  \tkzDefPoint(0,0){O}
  \tkzDefPoint(-2,0){A}
  \tkzDefPoint(2,0){B}
  \tkzDefPoint(0,2){C}

  \tkzDrawCircle(O,A)
  \tkzDrawSegments(A,B A,C B,C)
  \tkzMarkRightAngle[size=0.28](A,C,B)

  \tkzLabelPoints[left](A)
  \tkzLabelPoints[right](B)
  \tkzLabelPoints[above](C)
  \tkzLabelPoints[below](O)
\end{tikzpicture}
\end{center}

\key{Полезная связь.}
В прямоугольном треугольнике середина гипотенузы является центром
описанной окружности. Поэтому свойства медианы и свойства диаметра ---
два взгляда на одну и ту же конструкцию.

% =========================================================
\section*{\textcolor{darkaccent}{3. Если прямой угол разделён в отношении \(1:2\)}}
\sectionrule

Если некоторый луч делит прямой угол на части в отношении \(1:2\), то
\[
x+2x=90^\circ,
\]
откуда
\[
x=30^\circ.
\]

Следовательно, получаются углы
\[
30^\circ\quad\text{и}\quad60^\circ.
\]

Для прямоугольного треугольника с углами
\(30^\circ,60^\circ,90^\circ\):
\[
\boxed{
\text{катет напротив }30^\circ
:
\text{катет напротив }60^\circ
:
\text{гипотенуза}
=
1:\sqrt3:2
}.
\]

\textbf{Типовой пример.}
Медиана к гипотенузе равна \(m\) и делит прямой угол в отношении \(1:2\).

Так как медиана к гипотенузе равна половине гипотенузы,
\[
AB=2m.
\]
Углы равны \(30^\circ\) и \(60^\circ\), поэтому стороны треугольника:
\[
\boxed{m,\quad \sqrt3\,m,\quad 2m}.
\]

% =========================================================
\section*{\textcolor{darkaccent}{4. Теорема Пифагора и тригонометрия}}
\sectionrule

Для прямоугольного треугольника \(ABC\), где
\(\angle C=90^\circ\),
\[
\boxed{AB^2=AC^2+BC^2}.
\]

Если \(\alpha=\angle A\), то
\[
\sin\alpha=\frac{BC}{AB},
\qquad
\cos\alpha=\frac{AC}{AB},
\]
\[
\tg\alpha=\frac{BC}{AC},
\qquad
\ctg\alpha=\frac{AC}{BC}.
\]

\begin{center}
\renewcommand{\arraystretch}{1.3}
\begin{tabularx}{0.88\linewidth}{@{}lX@{}}
\toprule
\textbf{Функция} & \textbf{Что сравниваем} \\
\midrule
\(\sin\alpha\) &
противолежащий катет и гипотенузу \\

\(\cos\alpha\) &
прилежащий катет и гипотенузу \\

\(\tg\alpha\) &
противолежащий и прилежащий катеты \\

\(\ctg\alpha\) &
прилежащий и противолежащий катеты \\
\bottomrule
\end{tabularx}
\end{center}

\key{Самопроверка:}
для острого угла прямоугольного треугольника
\[
0<\sin\alpha<1,
\qquad
0<\cos\alpha<1.
\]

% =========================================================
\section*{\textcolor{darkaccent}{5. Площадь треугольника}}
\sectionrule

Основные формулы:
\[
\boxed{S=\frac12 ah_a},
\]
где \(h_a\) --- высота, проведённая к стороне \(a\).

Для прямоугольного треугольника с катетами \(a\) и \(b\):
\[
\boxed{S=\frac12 ab}.
\]

Если \(c\) --- гипотенуза, а \(h\) --- высота к гипотенузе, то
\[
\frac12 ab=\frac12 ch,
\]
поэтому
\[
\boxed{h=\frac{ab}{c}}.
\]

\subsection*{\textcolor{accent}{Формула Герона}}

Для треугольника со сторонами \(a,b,c\)
\[
p=\frac{a+b+c}{2},
\]
\[
\boxed{
S=\sqrt{p(p-a)(p-b)(p-c)}
}.
\]

Перед применением формулы необходимо проверить существование треугольника:
\[
a+b>c,\qquad
a+c>b,\qquad
b+c>a.
\]

% =========================================================
\section*{\textcolor{darkaccent}{6. Типовая задача: известны периметры двух частей}}
\sectionrule

Медиана прямоугольного треугольника, проведённая к гипотенузе,
разбивает его на два треугольника с периметрами \(8\) и \(9\).

Пусть
\[
CM=m.
\]
Тогда
\[
AM=BM=m,\qquad AB=2m.
\]

Обозначим катеты через \(x\) и \(y\). Из периметров:
\[
x+2m=8,
\qquad
y+2m=9.
\]
Следовательно,
\[
x=8-2m,
\qquad
y=9-2m.
\]

По теореме Пифагора:
\[
(8-2m)^2+(9-2m)^2=(2m)^2.
\]

После упрощения:
\[
(2m-5)(2m-29)=0.
\]

Получаем два алгебраических значения:
\[
m=\frac52
\qquad\text{или}\qquad
m=\frac{29}{2}.
\]

Второе значение невозможно, так как при нём катеты
\(8-2m\) и \(9-2m\) отрицательны. Поэтому
\[
m=\frac52.
\]

Тогда
\[
x=3,\qquad y=4,\qquad AB=5.
\]

\[
\boxed{\text{Стороны треугольника: }3,\ 4,\ 5.}
\]

\key{Важно.}
В геометрической задаче недостаточно получить корень уравнения.
Все длины должны быть положительными и согласованными с условием.

% =========================================================
\section*{\textcolor{darkaccent}{7. Медиана и высота в одном треугольнике}}
\sectionrule

Рассмотрим треугольник \(ABC\). К стороне \(AC\) проведены:

\begin{itemize}
  \item медиана \(BM\);
  \item высота \(BK\).
\end{itemize}

Пусть
\[
AM=BM,\qquad AB=1,\qquad BC=2.
\]

Так как \(BM\) --- медиана,
\[
AM=MC.
\]
По условию \(AM=BM\), следовательно
\[
AM=BM=CM.
\]

Точка \(M\) равноудалена от \(A,B,C\), поэтому \(AC\) является диаметром
описанной окружности и
\[
\angle ABC=90^\circ.
\]

По теореме Пифагора:
\[
AC=\sqrt{1^2+2^2}=\sqrt5.
\]

Следовательно,
\[
BM=\frac{AC}{2}=\frac{\sqrt5}{2}.
\]

Площадь:
\[
S_{ABC}
=
\frac12\cdot AB\cdot BC
=
1.
\]

С другой стороны,
\[
S_{ABC}
=
\frac12\cdot AC\cdot BK,
\]
поэтому
\[
\frac12\sqrt5\cdot BK=1,
\qquad
BK=\frac{2}{\sqrt5}.
\]

Треугольник \(BKM\) прямоугольный при \(K\), поэтому
\[
\cos\angle KBM
=
\frac{BK}{BM}
=
\frac{2/\sqrt5}{\sqrt5/2}
=
\boxed{\frac45}.
\]

% =========================================================
\section*{\textcolor{darkaccent}{8. Внешние квадраты на катетах}}
\sectionrule

Пусть \(ABC\) --- прямоугольный треугольник,
\[
\angle ACB=90^\circ,
\qquad
AC=a,\quad BC=b,\quad AB=c.
\]

На катетах \(AC\) и \(BC\) снаружи построены квадраты
\(ACDE\) и \(BCFG\).
Медиана \(CM\) продолжена до пересечения с прямой \(DF\) в точке \(N\).

\begin{center}
\begin{tikzpicture}[scale=0.72]
  \tkzDefPoint(0,0){C}
  \tkzDefPoint(0,4){A}
  \tkzDefPoint(1,0){B}

  \tkzDefPoint(-4,0){D}
  \tkzDefPoint(-4,4){E}

  \tkzDefPoint(0,-1){F}
  \tkzDefPoint(1,-1){G}

  \tkzDefMidPoint(A,B)
  \tkzGetPoint{M}

  \tkzInterLL(C,M)(D,F)
  \tkzGetPoint{N}

  \tkzDrawSegments(A,B B,C C,A)
  \tkzDrawPolygon(A,C,D,E)
  \tkzDrawPolygon(B,C,F,G)
  \tkzDrawSegment(D,F)
  \tkzDrawSegment(M,N)

  \tkzMarkRightAngle[size=0.25](A,C,B)
  \tkzMarkRightAngle[size=0.22](D,C,F)
  \tkzMarkRightAngle[size=0.22](D,N,C)

  \tkzLabelPoints[above](A,E)
  \tkzLabelPoints[right](B,G,M)
  \tkzLabelPoints[below](F,N)
  \tkzLabelPoints[left](D)
  \tkzLabelPoints[above left](C)
\end{tikzpicture}
\end{center}

Треугольники \(DCF\) и \(ACB\) имеют равные пары катетов:
\[
DC=AC=a,
\qquad
CF=CB=b,
\]
и прямые углы при \(C\). Поэтому
\[
\triangle DCF\cong\triangle ACB
\]
и
\[
DF=AB=c.
\]

Из угловых соотношений в равнобедренных треугольниках,
возникающих благодаря свойству медианы к гипотенузе,
получаем
\[
CN\perp DF.
\]

Следовательно, площадь \(\triangle DCF\) можно вычислить двумя способами:
\[
S_{DCF}
=
\frac12 ab
=
\frac12 DF\cdot CN.
\]

Так как \(DF=c\),
\[
\boxed{
CN=\frac{ab}{c}
}.
\]

По теореме Пифагора
\[
c=\sqrt{a^2+b^2},
\]
поэтому
\[
\boxed{
CN=\frac{ab}{\sqrt{a^2+b^2}}
}.
\]

Для катетов \(1\) и \(4\):
\[
CN=
\frac{1\cdot4}{\sqrt{1^2+4^2}}
=
\boxed{\frac{4}{\sqrt{17}}}
=
\boxed{\frac{4\sqrt{17}}{17}}.
\]

Заметим, что полученная величина совпадает с высотой прямоугольного
треугольника, проведённой к гипотенузе.

% =========================================================
\section*{\textcolor{darkaccent}{9. Универсальный алгоритм решения}}
\sectionrule

\begin{enumerate}[label=\textbf{\arabic*.}]
  \item \key{Отметьте конструкцию.}
  Найдите прямой угол, гипотенузу, медиану, высоту и середины сторон.

  \item \key{Переведите геометрию в равенства.}
  Например,
  \[
  M\text{ --- середина }AB
  \quad\Longrightarrow\quad
  AM=BM.
  \]

  \item \key{Проверьте, можно ли применить свойство медианы.}
  Только после установления прямоугольности допустимо писать
  \[
  CM=\frac{AB}{2}.
  \]

  \item \key{Выберите основной инструмент:}
  \[
  \text{Пифагор}
  \quad\text{или}\quad
  \text{тригонометрия}
  \quad\text{или}\quad
  \text{площадь}.
  \]

  \item Если в задаче известны три стороны общего треугольника,
  рассмотрите \key{формулу Герона}.

  \item Если получено уравнение, решите его и обязательно
  \key{проверьте геометрический смысл корней}.

  \item В конце выполните быструю проверку:
  длины положительны, гипотенуза --- наибольшая сторона,
  а найденные синус и косинус принадлежат интервалу \((0;1)\).
\end{enumerate}

% =========================================================
\section*{\textcolor{darkaccent}{10. Типичные ошибки}}
\sectionrule

\begin{itemize}
  \checkitem Писать
  \[
  CM=\frac{AB}{2}
  \]
  для произвольного треугольника. Это свойство относится к медиане,
  проведённой к \textbf{гипотенузе прямоугольного треугольника}.

  \checkitem Путать медиану и высоту:
  медиана делит сторону пополам, высота образует угол \(90^\circ\).

  \checkitem Из равенства \(AM=BM\) сразу делать вывод о прямом угле,
  не используя дополнительно тот факт, что \(M\) --- середина стороны.

  \checkitem При отношении углов \(1:2\) забывать, что их сумма равна
  именно \(90^\circ\), если они составляют прямой угол.

  \checkitem В формулах \(\sin\) и \(\cos\) выбирать катет
  безотносительно рассматриваемого угла.

  \checkitem При вычислении площади использовать высоту,
  проведённую не к той стороне, которая выбрана основанием.

  \checkitem Сохранять алгебраический корень, дающий отрицательную
  или нулевую длину.

  \checkitem Применять формулу Герона, не проверив неравенства треугольника.

  \checkitem В сложной конфигурации считать отрезки равными
  только по внешнему виду рисунка, без доказанного свойства.
\end{itemize}

% =========================================================
% TRAINING
% =========================================================
\clearpage
\section*{\textcolor{darkaccent}{Тренировочный блок}}
\sectionrule

Выполните задания без подсказок. По возможности сначала сделайте
схематический чертёж и отметьте на нём все известные равенства.

\begin{enumerate}[label=\textbf{\arabic*.}]

\item[]
\textcolor{accent}{\textbf{Средний уровень}}

\item
В прямоугольном треугольнике гипотенуза равна \(14\).
Найдите медиану, проведённую к гипотенузе.

\item
Медиана, проведённая к гипотенузе прямоугольного треугольника,
равна \(5\) и делит прямой угол в отношении \(1:2\).
Найдите все стороны треугольника.

\item
Катеты прямоугольного треугольника равны \(6\) и \(8\).
Найдите:
\[
\text{а) гипотенузу;}
\qquad
\text{б) медиану к гипотенузе;}
\]
\[
\text{в) синус меньшего острого угла.}
\]

\item[]
\textcolor{accent}{\textbf{Выше среднего}}

\item
Медиана прямоугольного треугольника, проведённая к гипотенузе,
разбивает его на два треугольника с периметрами \(16\) и \(18\).
Найдите стороны исходного треугольника.

\item
В треугольнике \(ABC\) к стороне \(AC\) проведены высота \(BK\)
и медиана \(BM\). Известно, что
\[
AM=BM,\qquad AB=6,\qquad BC=8.
\]
Найдите
\[
\cos\angle KBM.
\]

\item
Точки \(A\) и \(B\) --- концы диаметра окружности,
\[
AB=10.
\]
Точка \(C\) лежит на окружности, причём
\[
AC=6.
\]
Найдите \(BC\) и медиану \(CM\), проведённую к стороне \(AB\).

\item
Стороны треугольника равны
\[
13,\qquad14,\qquad15.
\]
Найдите площадь треугольника и высоту, проведённую к стороне \(14\).

\item
Один из острых углов прямоугольного треугольника равен \(30^\circ\),
а площадь треугольника равна
\[
9\sqrt3.
\]
Найдите все его стороны.

\item
Вне прямоугольного треугольника \(ABC\) на катетах \(AC\) и \(BC\)
построены квадраты \(ACDE\) и \(BCFG\).
Продолжение медианы \(CM\) пересекает прямую \(DF\) в точке \(N\).
Найдите \(CN\), если катеты исходного треугольника равны
\[
5\quad\text{и}\quad12.
\]

\item[]
\textcolor{accent}{\textbf{Сложный уровень}}

\item
Медиана \(CM\) прямоугольного треугольника \(ABC\),
проведённая к гипотенузе, разбивает его на два треугольника
с периметрами \(16\) и \(18\).

Вне треугольника на катетах \(AC\) и \(BC\) построены квадраты
\(ACDE\) и \(BCFG\).
Продолжение медианы \(CM\) пересекает прямую \(DF\) в точке \(N\).

Найдите длину \(CN\).

\end{enumerate}

% =========================================================
% ANSWERS
% =========================================================
\clearpage
\section*{\textcolor{darkaccent}{Ответы к тренировочному блоку}}
\sectionrule

\begin{table}[ht]
\centering
\caption{Ответы для самопроверки}
\renewcommand{\arraystretch}{1.45}
\begin{tabularx}{\linewidth}{@{}cX@{}}
\toprule
\textbf{№} & \textbf{Ответ} \\
\midrule

1 &
\[
7
\]
\\

2 &
\[
5,\qquad5\sqrt3,\qquad10
\]
\\

3 &
\[
\text{а) }10;
\qquad
\text{б) }5;
\qquad
\text{в) }\frac35
\]
\\

4 &
\[
6,\qquad8,\qquad10
\]
\\

5 &
\[
\frac{24}{25}
\]
\\

6 &
\[
BC=8,\qquad CM=5
\]
\\

7 &
\[
S=84,\qquad h=12
\]
\\

8 &
\[
3\sqrt2,\qquad3\sqrt6,\qquad6\sqrt2
\]
\\

9 &
\[
CN=\frac{60}{13}
\]
\\

10 &
\[
CN=\frac{24}{5}
\]
\\

\bottomrule
\end{tabularx}
\end{table}

\vfill

\begin{center}
{\color{accent!65}\rule{0.62\linewidth}{0.8pt}}

\vspace{0.5em}

\small
Перед проверкой ответа убедитесь, что использованное свойство
действительно применимо к данной конфигурации.
\end{center}

\end{document}